\section{Tactile Observations, Sensors, and Data Regimes}
\label{sec:sensing}

\subsection{The sensor is part of the representation problem}

Tactile learning spans RGB or marker-based deformation images, elastomer flow,
taxel pressure arrays, event-based tactile cameras, distributed fingertip force,
whole-hand skins, and wrist or interface wrench. These signals differ not only in
dimension but in invariances, latency, hysteresis, saturation, calibration, and
spatial support. A ``tactile token'' is therefore not a sensor-neutral primitive.
Cross-sensor learning must decide which information is shared---contact geometry,
normal/shear deformation, dynamics, semantic material cues---and which remains
sensor- or embodiment-specific.

Hardware lifetime is part of this invariance problem. A replaceable-cartridge
vision-based fingertip reports markedly improved resistance to abrasion and
concentrated loading under its accelerated tests~\cite{cottrell2026durable}, while
a uniformly illuminated visuo-tactile sensor couples depth reconstruction with
spatiotemporal contact-state and slip-direction classification
~\cite{ma2026slip}. Durability and illumination stability do not by themselves
solve representation transfer, but degradation state, cartridge identity, and
lighting/calibration changes should be recorded as sensor-domain variables rather
than absorbed into unexplained model error.

Vision-based tactile sensing has enabled large-scale self-supervision because a
frame can be processed with mature visual backbones. Sparsh learns general touch
representations from large unlabeled tactile collections~\cite{higuera2025sparsh};
T-Dex uses robotic play to pretrain features useful for dexterous
control~\cite{guzey2023tdex}; UniTouch and AnyTouch align heterogeneous sensors and
static/dynamic content~\cite{yang2024unitouch,feng2025anytouch}; and T3 studies
transferable tactile Transformers~\cite{zhao2024t3}. These works make local tactile
content more reusable, but most cross-sensor objectives do not preserve a global
hand--object--hand contact identity by themselves.

Distributed and whole-hand sensors shift the difficulty. A fingertip array has a
natural patch boundary; a skin spanning phalanges and palm requires morphology
and sensor-layout metadata. HTT explicitly addresses heterogeneous tactile
sensors~\cite{bi2026htt}; HT-Bench expands full-hand tactile representation with
egocentric vision~\cite{huang2026htbench}; FTP-1 introduces morphology-aware
tactile tokens across sensors and embodiments~\cite{yuan2026ftp1}. Such work is
directly relevant to bimanual transfer because the two hands may have different
coverage, calibration, or missing channels.

\subsection{Force estimates, measured wrench, and ground truth}

Three quantities are frequently conflated:

\begin{itemize}
  \item a raw tactile deformation or taxel field;
  \item a wrench estimated from that tactile signal by calibration or a learned
  model; and
  \item an independently measured interface wrench from a dedicated force/torque
  sensor or instrumented object.
\end{itemize}

The distinction determines what can serve as a target. Training a model to predict
an estimate derived from the same tactile image measures compression or
consistency, not independent physical validity. Likewise, a simulator's net
fingertip contact force is useful supervision but is not equivalent to a real
tactile array. PhysGraph, for example, studies graph-structured bimanual policy
learning with simulator state and fingertip resultant forces rather than raw
tactile images~\cite{li2026physgraph}. This makes it relevant to relational control
while limiting conclusions about tactile representation transfer.

Force inference can also drift even when the contact geometry appears unchanged.
Relaxation-aware soft-gripper sensing uses onboard vision, infrared thermography,
and a temperature-coupled viscoelastic model to estimate force during sustained
holds~\cite{wang2026relaxation}. It shows why contact duration, temperature, and
loading history belong in the observation audit. An instantaneous tactile or
force frame may be insufficient for a long hold, and agreement with a same-frame
proxy should not be mistaken for relaxation-robust force estimation.

For load-transfer research, the strongest target hierarchy is:

\begin{enumerate}
  \item independent instrumented-object or interface force/torque where feasible;
  \item calibrated robot wrist/hand wrench with explicit coordinate transforms and
  gravity compensation;
  \item simulator wrench for pretraining or privileged-to-observation
  distillation, followed by a real validation subset;
  \item proxy load share or contact phase, with the proxy's construction and
  circularity documented.
\end{enumerate}

\subsection{Data regimes and their affordances}

Human-collected datasets can provide semantic and visual diversity at scale.
Touch and Go aligns human-collected touch and vision~\cite{yang2022touchandgo}; OAKINK2
captures complex bimanual hands--object activity~\cite{zhan2024oakink2}; and TACO
focuses bimanual tool--action--object understanding~\cite{liu2024taco}. Their
strength is breadth, while precise synchronized robot proprioception, calibrated
loads, and intervention labels may be absent.

Robot-collected data provides action, embodiment, and timing. VTDexManip offers a
visual-tactile pretraining and dexterous manipulation benchmark that includes
bimanual handover~\cite{liu2025vtdexmanip}. T-Rex reports a large tactile-reactive
robot dataset and bimanual capabilities~\cite{niu2026trex}. RCT targets
robot-collected touch--vision--language generalization~\cite{he2026rct}, while
TactiDex and HT-Bench expand real, whole-hand tactile evaluation
~\cite{ni2026tactidex,huang2026htbench}. As of the artifact audit on 8 August
2026, the TactiDex author project page stated ACM MM 2026 acceptance, but no
downloadable dataset, code, or model artifact had been confirmed; it should not
be assumed to be an executable first-month substrate. Because several of these appeared as
preprints or accepted works near the corpus freeze, exact public availability,
license, and final metadata must be checked before they are treated as
reproduction substrates.

Two new datasets illustrate complementary physical labels. ViTacPhys uses human
visual--tactile demonstrations to estimate mass class, friction class, and
stiffness before conditioning an adaptive grasping policy
~\cite{liu2026vitacphys}. SoftVTBench synchronizes policy-visible visual--tactile
signals with evaluator-only finite-element state and defines a deformation-aware
success criterion~\cite{jing2026softvtbench}. The former makes physical-property
probes concrete; the latter shows that task completion can coexist with excessive
deformation. Artifact status remains part of the evidence: at the 30 August audit,
ViTacPhys code/data were still listed as forthcoming, and the public SoftVTBench
release described a smaller dataset version than the current manuscript.

Simulation is essential for contact labels and topology interventions. Tactile
Genesis targets scalable tactile simulation across layouts~\cite{chung2026tactilegenesis};
generic graph simulators demonstrate how object/particle relations can support
long-horizon physical prediction~\cite{sanchez2020gns,allen2022contact}. Simulation
can generate edge, force, slip, and counterfactual labels that are prohibitively
expensive in the real world. Its risk is an unobserved shortcut: a model may use
precise object state, collision geometry, contact impulses, or future reference
that will not exist at deployment.

\begin{table}[t]
\centering
\caption{Representative data regimes and what they can and cannot establish for
cross-interface representation. Status reflects the corpus freeze.}
\label{tab:data_regimes}
\begin{tabularx}{\linewidth}{P{0.18\linewidth} P{0.18\linewidth} P{0.20\linewidth} Y}
\toprule
\textbf{Resource / line} & \textbf{Regime} & \textbf{Main strength} & \textbf{Caution for bimanual load-path claims} \\
\midrule
Touch and Go~\cite{yang2022touchandgo} & human-collected vision--touch & broad multimodal diversity & limited robot action, proprioception, and calibrated cross-hand load labels \\
T-Dex / Sparsh~\cite{guzey2023tdex,higuera2025sparsh} & robot touch / large unlabeled tactile data & reusable local encoders & downstream topology not specified by pretraining objective \\
VTDexManip~\cite{liu2025vtdexmanip} & simulated visual--tactile benchmark & reproducible pretraining and handover task & sim observation and contact fidelity must match intended deployment \\
T-Rex~\cite{niu2026trex} & large real robot tactile data & temporally rich bimanual control evidence & dataset subset, synchronization, and exact labels require audit \\
OAKINK2 / TACO~\cite{zhan2024oakink2,liu2024taco} & bimanual human/object datasets & topology and semantic task breadth & tactile load evidence is not their primary signal \\
HT-Bench / TactiDex~\cite{huang2026htbench,ni2026tactidex} & whole-hand real tactile benchmarks & coverage beyond fingertips & recent status and resource release should be verified per artifact \\
ViTacPhys~\cite{liu2026vitacphys} & human visual--tactile demonstrations and robot transfer & explicit mass, friction, and stiffness probes & semantic priors and unreleased artifacts must be separated from tactile evidence \\
SoftVTBench~\cite{jing2026softvtbench} & deformable-object visual--tactile benchmark & independent deformation state and safety-aware success & public release version must be matched to the reported manuscript scale \\
Tactile Genesis~\cite{chung2026tactilegenesis} & scalable tactile simulation & layout variation and dense labels & sim-to-real contact appearance and dynamics gap \\
\bottomrule
\end{tabularx}
\end{table}

\subsection{Synchronization is not a preprocessing detail}

Bimanual tactile data is intrinsically multi-rate. Cameras, tactile arrays,
joint encoders, wrist wrenches, and policy actions may run at different rates and
have different transport delays. A 50--100 ms misalignment can change an apparent
causal relation at contact onset or release. Dataset construction should therefore
retain raw timestamps, clock source, interpolation policy, sensor-quality masks,
and commanded versus measured action times. Synchronization should be evaluated
with event-based sanity checks---for example, whether tactile onset precedes an
object acceleration that it could physically cause---rather than only by visually
aligned plots.

The split unit is equally important. Frame-level random splits leak object,
episode, grasp, and contact-sequence identity into the test set. Representation
claims require episode-level splits at minimum, with object, session, operator,
sensor, embodiment, and contact-sequence holds-outs selected according to the
claim. For cross-topology transfer, the complete topology family---not just a
motion instance---must be absent from training.

\subsection{A minimum interoperable schema}

A useful cross-dataset schema should keep (i) raw tactile observations and
quality masks; (ii) calibrated per-sensor metadata; (iii) joint states and link
poses; (iv) object/tool/workpiece identifiers and optional perception-derived
  poses; (v) actions and timestamps; (vi) phase, interface, slip, and wrench labels
  with provenance; (vii) temperature, contact duration, maintenance/calibration
  state, and material metadata when relaxation or wear matters; and (viii) split
  membership. Coordinate frames must be explicit.
Without this information, a study can compare end-to-end policies, but it cannot
reliably compare structured interface representations.
